\documentclass[11pt]{article}
\usepackage[margin=1in]{geometry}
\usepackage{amsmath,amssymb}
\usepackage[T1]{fontenc}
\usepackage{lmodern}
\usepackage{microtype}
\usepackage[hidelinks]{hyperref}

\title{Literature-Implied Partial Answer for arXiv:1309.4684}
\author{agent\_lane\_07}
\date{2026-06-20}

\begin{document}
\maketitle

\section*{Status}

\textbf{literature\_implied\_answer (partial subcase).}
This note records a later literature theorem which answers one named family
from Bendova's final open-ended question. It is not claimed as a new result of
this run.

\section*{Source Question}

The source is Hana Bendova, \emph{Quantitative Grothendieck property},
arXiv:1309.4684. In Section 4, after proving that
\(\ell_\infty(\Gamma)\), sigma-finite \(L^\infty(\mu)\), and quotients of
\(\ell_\infty(\Gamma)\) are \(1\)-Grothendieck, Bendova asks whether the
other Grothendieck spaces mentioned in the introduction have the quantitative
version as well. The introduction names, among others, \(C(K)\) for compact
\(F\)-spaces, weak \(L^p\) spaces, \(H^\infty\), and Haydon's Grothendieck
space not containing \(\ell_\infty\).

\section*{Supporting Literature}

The supporting paper is Jindrich Lechner, \emph{1-Grothendieck \(C(K)\)
spaces}, arXiv:1511.02202. Its main theorem states:
\[
  \text{if } S \text{ is a Haydon space, then } C(S) \text{ is }
  1\text{-Grothendieck.}
\]
Lechner defines a Haydon space as a totally disconnected compact space whose
algebra of clopen subsets has the subsequential completeness property. The
paper also points out that Haydon had constructed such a space \(K\) with
\(C(K)\) not containing an isomorphic copy of \(\ell_\infty\). Therefore the
specific Haydon example named in Bendova's introduction has an affirmative
quantitative answer, with constant \(1\).

Lechner further records that \(C(K)\) is \(1\)-Grothendieck whenever \(K\) is
sigma-Stonean, and in particular whenever \(K\) is extremally disconnected.
These are additional \(C(K)\)-side positive cases, but they do not by
themselves cover all compact \(F\)-spaces.

\section*{Scope and Remaining Cases}

The identification gives the following partial answer to arXiv:1309.4684:
\[
  \boxed{\text{Haydon-space } C(K) \text{ examples are } 1
  \text{-Grothendieck}.}
\]
In particular, Haydon's Grothendieck space without \(\ell_\infty\) satisfies
Bendova's quantitative Grothendieck property with the optimal constant allowed
by the definition.

The packet does not settle the other families listed by Bendova: general
compact \(F\)-space \(C(K)\) spaces, weak \(L^p\) spaces, or \(H^\infty\).
Lechner explicitly leaves the other examples from his introduction open after
the Haydon-space theorem. During this bounded pass, local registry searches and
web searches for close phrases around \(1\)-Grothendieck, quantitative
Grothendieck, \(H^\infty\), weak \(L^p\), compact \(F\)-spaces, and Haydon
spaces found Lechner's theorem as the decisive supporting result and did not
locate a complete later answer for the remaining Bendova examples.

\section*{References}

\begin{enumerate}
\item H. Bendova, \emph{Quantitative Grothendieck property},
  arXiv:1309.4684.
\item J. Lechner, \emph{1-Grothendieck \(C(K)\) spaces},
  arXiv:1511.02202.
\end{enumerate}

\end{document}
